%%%%%%%%%%%%%%%%%%%%%%%%%%%%%%%%%%%%%%%%%%%%%%%%%%%%%%%%%%%%
% 14.16 MATHEMATICAL TERMINOLOGY
% The Composition of Advanced Mathematics
%%%%%%%%%%%%%%%%%%%%%%%%%%%%%%%%%%%%%%%%%%%%%%%%%%%%%%%%%%%%

\section{Mathematical Terminology}
\label{sec:mathematical-terminology}

\prerequisite{
Basic mathematical notation and familiarity with the major
branches of mathematics.
}

\learningobjectives{
By the end of this reference section, the reader should be able to:
\begin{itemize}
    \item recognize common mathematical vocabulary;
    \item distinguish closely related terms;
    \item identify terminology associated with algebra, geometry,
          analysis, probability, statistics, and applied mathematics;
    \item interpret common words used in theorem statements and proofs;
    \item recognize standard computational and modeling terminology;
    \item use mathematical terminology more precisely in written solutions.
\end{itemize}
}

%%%%%%%%%%%%%%%%%%%%%%%%%%%%%%%%%%%%%%%%%%%%%%%%%%%%%%%%%%%%
\subsection{Purpose of Mathematical Terminology}
\label{subsec:terminology-purpose}
%%%%%%%%%%%%%%%%%%%%%%%%%%%%%%%%%%%%%%%%%%%%%%%%%%%%%%%%%%%%

Mathematics depends heavily on precise vocabulary.

Words such as

\[
\boxed{
\text{function},
\quad
\text{continuous},
\quad
\text{independent},
\quad
\text{normal},
\quad
\text{regular}
}
\]

have specific technical meanings that may differ from ordinary language.

This section groups common terminology by mathematical subject.

The alphabetical Glossary in Section~14.17 provides the broader
term-by-term reference.

%%%%%%%%%%%%%%%%%%%%%%%%%%%%%%%%%%%%%%%%%%%%%%%%%%%%%%%%%%%%
\subsection{Logical Terminology}
\label{subsec:logical-terminology-reference}
%%%%%%%%%%%%%%%%%%%%%%%%%%%%%%%%%%%%%%%%%%%%%%%%%%%%%%%%%%%%

\textbf{Statement.}
A declarative sentence that is either true or false.

\medskip

\textbf{Proposition.}
A statement asserted and usually proved.

\medskip

\textbf{Predicate.}
A statement containing one or more variables whose truth depends on those
variables.

\medskip

\textbf{Hypothesis.}
An assumption made in a theorem, implication, or argument.

\medskip

\textbf{Conclusion.}
The statement claimed to follow from the hypotheses.

\medskip

\textbf{Implication.}
A logical statement of form

\[
\boxed{
P\Rightarrow Q.
}
\]

\medskip

\textbf{Converse.}
The converse of

\[
P\Rightarrow Q
\]

is

\[
\boxed{
Q\Rightarrow P.
}
\]

\medskip

\textbf{Contrapositive.}
The contrapositive is

\[
\boxed{
\neg Q\Rightarrow\neg P.
}
\]

It is logically equivalent to the original implication.

\medskip

\textbf{Contradiction.}
A statement that cannot be true, often used as the endpoint of proof by
contradiction.

\medskip

\textbf{Necessary condition.}
\(P\) is necessary for \(Q\) when

\[
\boxed{
Q\Rightarrow P.
}
\]

\medskip

\textbf{Sufficient condition.}
\(P\) is sufficient for \(Q\) when

\[
\boxed{
P\Rightarrow Q.
}
\]

\medskip

\textbf{Necessary and sufficient.}
Both implications hold:

\[
\boxed{
P\iff Q.
}
\]

%%%%%%%%%%%%%%%%%%%%%%%%%%%%%%%%%%%%%%%%%%%%%%%%%%%%%%%%%%%%
\subsection{Proof Terminology}
\label{subsec:proof-terminology-reference}
%%%%%%%%%%%%%%%%%%%%%%%%%%%%%%%%%%%%%%%%%%%%%%%%%%%%%%%%%%%%

\textbf{Theorem.}
An important mathematical statement established by proof.

\medskip

\textbf{Lemma.}
A supporting theorem usually used to prove another result.

\medskip

\textbf{Corollary.}
A result that follows relatively directly from a theorem.

\medskip

\textbf{Conjecture.}
A mathematical statement believed to be true but not yet proved.

\medskip

\textbf{Counterexample.}
An example showing that a universal statement is false.

\medskip

\textbf{Direct proof.}
A proof that starts from the hypotheses and derives the conclusion.

\medskip

\textbf{Proof by contrapositive.}
A proof of

\[
P\Rightarrow Q
\]

by proving

\[
\neg Q\Rightarrow\neg P.
\]

\medskip

\textbf{Proof by contradiction.}
A proof that assumes the negation of the desired conclusion and derives an
impossibility.

\medskip

\textbf{Mathematical induction.}
A proof method involving a base case and an inductive step.

\medskip

\textbf{Invariant.}
A quantity or property that remains unchanged under a specified process or
transformation.

%%%%%%%%%%%%%%%%%%%%%%%%%%%%%%%%%%%%%%%%%%%%%%%%%%%%%%%%%%%%
\subsection{Set-Theoretic Terminology}
\label{subsec:set-theoretic-terminology}
%%%%%%%%%%%%%%%%%%%%%%%%%%%%%%%%%%%%%%%%%%%%%%%%%%%%%%%%%%%%

\textbf{Set.}
A collection of distinct mathematical objects.

\medskip

\textbf{Element.}
An object belonging to a set.

\[
\boxed{
x\in A
}
\]

means \(x\) is an element of \(A\).

\medskip

\textbf{Subset.}

\[
\boxed{
A\subseteq B
}
\]

means every element of \(A\) belongs to \(B\).

\medskip

\textbf{Proper subset.}
A subset that is not equal to the containing set.

\medskip

\textbf{Union.}

\[
\boxed{
A\cup B
}
\]

contains elements belonging to \(A\), \(B\), or both.

\medskip

\textbf{Intersection.}

\[
\boxed{
A\cap B
}
\]

contains elements common to both sets.

\medskip

\textbf{Complement.}
Elements in the ambient universe that are not in the specified set.

\medskip

\textbf{Difference.}

\[
\boxed{
A\setminus B
}
\]

contains elements of \(A\) that are not in \(B\).

\medskip

\textbf{Power set.}

\[
\boxed{
\mathcal{P}(A)
}
\]

is the set of all subsets of \(A\).

\medskip

\textbf{Cartesian product.}

\[
\boxed{
A\times B
}
\]

is the set of ordered pairs

\[
(a,b)
\]

with \(a\in A\) and \(b\in B\).

\medskip

\textbf{Cardinality.}
The size of a set, written for finite sets as

\[
\boxed{
|A|.
}
\]

%%%%%%%%%%%%%%%%%%%%%%%%%%%%%%%%%%%%%%%%%%%%%%%%%%%%%%%%%%%%
\subsection{Function Terminology}
\label{subsec:function-terminology-reference}
%%%%%%%%%%%%%%%%%%%%%%%%%%%%%%%%%%%%%%%%%%%%%%%%%%%%%%%%%%%%

\textbf{Function.}
A mapping assigning each element of a domain exactly one element of a
codomain.

\[
\boxed{
f:A\to B.
}
\]

\medskip

\textbf{Domain.}
The set of allowed inputs.

\medskip

\textbf{Codomain.}
The declared target set of a function.

\medskip

\textbf{Range or image.}
The values actually attained:

\[
\boxed{
f(A).
}
\]

\medskip

\textbf{Preimage.}
For \(S\subseteq B\),

\[
\boxed{
f^{-1}(S)
}
\]

is the set of inputs mapped into \(S\).

\medskip

\textbf{Injective.}
Distinct inputs have distinct outputs.

Also called one-to-one.

\medskip

\textbf{Surjective.}
Every element of the codomain is attained.

Also called onto.

\medskip

\textbf{Bijective.}
Both injective and surjective.

A bijection has an inverse function.

\medskip

\textbf{Composition.}

\[
\boxed{
(f\circ g)(x)
=
f(g(x)).
}
\]

\medskip

\textbf{Identity function.}

\[
\boxed{
\operatorname{id}_A(x)=x.
}
\]

%%%%%%%%%%%%%%%%%%%%%%%%%%%%%%%%%%%%%%%%%%%%%%%%%%%%%%%%%%%%
\subsection{Number Terminology}
\label{subsec:number-terminology-reference}
%%%%%%%%%%%%%%%%%%%%%%%%%%%%%%%%%%%%%%%%%%%%%%%%%%%%%%%%%%%%

\textbf{Natural number.}
A counting number, with the convention concerning inclusion of \(0\)
specified by context.

\medskip

\textbf{Integer.}
An element of

\[
\boxed{
\mathbb{Z}.
}
\]

\medskip

\textbf{Rational number.}
A number expressible as

\[
\boxed{
\frac pq,
\qquad
p,q\in\mathbb{Z},
\quad
q\neq0.
}
\]

\medskip

\textbf{Irrational number.}
A real number that is not rational.

\medskip

\textbf{Real number.}
An element of

\[
\boxed{
\mathbb{R}.
}
\]

\medskip

\textbf{Complex number.}
A number of form

\[
\boxed{
a+bi,
\qquad
i^2=-1.
}
\]

\medskip

\textbf{Prime number.}
An integer greater than \(1\) with exactly two positive divisors.

\medskip

\textbf{Composite number.}
An integer greater than \(1\) that is not prime.

\medskip

\textbf{Algebraic number.}
A complex number satisfying a nonzero polynomial equation with rational
coefficients.

\medskip

\textbf{Transcendental number.}
A complex number that is not algebraic.

%%%%%%%%%%%%%%%%%%%%%%%%%%%%%%%%%%%%%%%%%%%%%%%%%%%%%%%%%%%%
\subsection{Algebraic Terminology}
\label{subsec:algebraic-terminology-reference}
%%%%%%%%%%%%%%%%%%%%%%%%%%%%%%%%%%%%%%%%%%%%%%%%%%%%%%%%%%%%

\textbf{Variable.}
A symbol representing a value that may vary.

\medskip

\textbf{Constant.}
A quantity treated as fixed in the relevant context.

\medskip

\textbf{Coefficient.}
A multiplicative factor attached to a term.

\medskip

\textbf{Term.}
A component of an algebraic sum.

\medskip

\textbf{Expression.}
A mathematically meaningful combination of quantities and operations.

\medskip

\textbf{Equation.}
A statement asserting equality of two expressions.

\medskip

\textbf{Identity.}
An equation true for every value in its stated domain.

\medskip

\textbf{Solution.}
A value or object satisfying an equation or problem condition.

\medskip

\textbf{Root or zero.}
A value \(r\) satisfying

\[
\boxed{
f(r)=0.
}
\]

\medskip

\textbf{Polynomial.}
An expression

\[
\boxed{
a_nx^n+\cdots+a_1x+a_0
}
\]

with nonnegative integer exponents.

\medskip

\textbf{Degree.}
The largest exponent with nonzero coefficient in a polynomial.

%%%%%%%%%%%%%%%%%%%%%%%%%%%%%%%%%%%%%%%%%%%%%%%%%%%%%%%%%%%%
\subsection{Abstract Algebra Terminology}
\label{subsec:abstract-algebra-terminology}
%%%%%%%%%%%%%%%%%%%%%%%%%%%%%%%%%%%%%%%%%%%%%%%%%%%%%%%%%%%%

\textbf{Binary operation.}
A rule combining two elements of a set to produce another element of the
same set.

\medskip

\textbf{Group.}
A set with an associative binary operation, identity, and inverses.

\medskip

\textbf{Abelian group.}
A group whose operation is commutative.

\medskip

\textbf{Subgroup.}
A subset that forms a group under the inherited operation.

\medskip

\textbf{Ring.}
An algebraic structure with addition and multiplication satisfying the ring
axioms.

\medskip

\textbf{Integral domain.}
A commutative ring with identity and no zero divisors.

\medskip

\textbf{Field.}
A commutative ring in which every nonzero element has a multiplicative
inverse.

\medskip

\textbf{Homomorphism.}
A map preserving the relevant algebraic operations.

\medskip

\textbf{Isomorphism.}
A bijective homomorphism whose inverse preserves the structure.

\medskip

\textbf{Kernel.}
Elements mapped to the identity or zero under a homomorphism.

\medskip

\textbf{Quotient structure.}
A structure formed by identifying elements according to an equivalence
relation compatible with the operations.

%%%%%%%%%%%%%%%%%%%%%%%%%%%%%%%%%%%%%%%%%%%%%%%%%%%%%%%%%%%%
\subsection{Linear Algebra Terminology}
\label{subsec:linear-algebra-terminology}
%%%%%%%%%%%%%%%%%%%%%%%%%%%%%%%%%%%%%%%%%%%%%%%%%%%%%%%%%%%%

\textbf{Vector.}
An element of a vector space.

\medskip

\textbf{Scalar.}
An element of the field over which a vector space is defined.

\medskip

\textbf{Vector space.}
A set equipped with vector addition and scalar multiplication satisfying the
vector-space axioms.

\medskip

\textbf{Linear combination.}

\[
\boxed{
c_1v_1+\cdots+c_nv_n.
}
\]

\medskip

\textbf{Span.}
The set of all linear combinations of a collection of vectors.

\medskip

\textbf{Linearly independent.}
Vectors \(v_1,\ldots,v_n\) are linearly independent if

\[
c_1v_1+\cdots+c_nv_n=0
\]

implies

\[
\boxed{
c_1=\cdots=c_n=0.
}
\]

\medskip

\textbf{Basis.}
A linearly independent spanning set.

\medskip

\textbf{Dimension.}
The number of vectors in a basis of a finite-dimensional vector space.

\medskip

\textbf{Matrix rank.}
The dimension of the column space or row space.

\medskip

\textbf{Null space.}
The set

\[
\boxed{
\{x:Ax=0\}.
}
\]

\medskip

\textbf{Eigenvalue.}
A scalar \(\lambda\) satisfying

\[
\boxed{
Av=\lambda v
}
\]

for some nonzero vector \(v\).

\medskip

\textbf{Eigenvector.}
A nonzero vector satisfying the eigenvalue equation.

\medskip

\textbf{Diagonalizable.}
A matrix similar to a diagonal matrix.

\medskip

\textbf{Orthogonal.}
Vectors \(u,v\) are orthogonal if

\[
\boxed{
u\cdot v=0.
}
\]

%%%%%%%%%%%%%%%%%%%%%%%%%%%%%%%%%%%%%%%%%%%%%%%%%%%%%%%%%%%%
\subsection{Geometry Terminology}
\label{subsec:geometry-terminology-reference}
%%%%%%%%%%%%%%%%%%%%%%%%%%%%%%%%%%%%%%%%%%%%%%%%%%%%%%%%%%%%

\textbf{Point.}
A location with no intrinsic size.

\medskip

\textbf{Line.}
A one-dimensional straight geometric object.

\medskip

\textbf{Plane.}
A flat two-dimensional geometric object.

\medskip

\textbf{Angle.}
A measure of rotation between rays or directions.

\medskip

\textbf{Congruent.}
Geometric figures having the same size and shape.

\medskip

\textbf{Similar.}
Figures having the same shape up to scaling.

\medskip

\textbf{Parallel.}
Objects lying in directions that do not intersect in Euclidean geometry.

\medskip

\textbf{Perpendicular.}
Objects meeting at a right angle.

\medskip

\textbf{Convex set.}
A set \(C\) for which

\[
\boxed{
(1-t)x+ty\in C
}
\]

whenever \(x,y\in C\) and \(0\leq t\leq1\).

\medskip

\textbf{Geodesic.}
A locally distance-minimizing or straightest path in a geometric space,
depending on context.

%%%%%%%%%%%%%%%%%%%%%%%%%%%%%%%%%%%%%%%%%%%%%%%%%%%%%%%%%%%%
\subsection{Topology Terminology}
\label{subsec:topology-terminology-reference}
%%%%%%%%%%%%%%%%%%%%%%%%%%%%%%%%%%%%%%%%%%%%%%%%%%%%%%%%%%%%

\textbf{Open set.}
A set belonging to the topology of a topological space.

\medskip

\textbf{Closed set.}
A set whose complement is open.

\medskip

\textbf{Neighborhood.}
A set containing an open set around a specified point.

\medskip

\textbf{Interior.}
The largest open set contained in a given set.

\medskip

\textbf{Closure.}
The smallest closed set containing a given set.

\medskip

\textbf{Boundary.}
Points that can be approached from both the set and its complement.

\medskip

\textbf{Connected.}
A space that cannot be separated into two disjoint nonempty open pieces.

\medskip

\textbf{Path connected.}
Any two points can be joined by a continuous path.

\medskip

\textbf{Compact.}
A space for which every open cover has a finite subcover.

\medskip

\textbf{Homeomorphism.}
A bijective continuous map with continuous inverse.

Two homeomorphic spaces are topologically equivalent.

%%%%%%%%%%%%%%%%%%%%%%%%%%%%%%%%%%%%%%%%%%%%%%%%%%%%%%%%%%%%
\subsection{Limit and Continuity Terminology}
\label{subsec:limit-continuity-terminology}
%%%%%%%%%%%%%%%%%%%%%%%%%%%%%%%%%%%%%%%%%%%%%%%%%%%%%%%%%%%%

\textbf{Limit.}
The value approached by a function or sequence in a specified limiting
process.

\medskip

\textbf{One-sided limit.}
A limit obtained by approaching from one direction.

\medskip

\textbf{Continuous.}
A function \(f\) is continuous at \(a\) when

\[
\boxed{
\lim_{x\to a}f(x)=f(a).
}
\]

\medskip

\textbf{Uniformly continuous.}
A stronger form of continuity in which one \(\delta\) works uniformly over
the domain for each \(\varepsilon\).

\medskip

\textbf{Discontinuity.}
A point where continuity fails.

\medskip

\textbf{Asymptote.}
A line or other reference curve approached by a function in a specified
limiting sense.

%%%%%%%%%%%%%%%%%%%%%%%%%%%%%%%%%%%%%%%%%%%%%%%%%%%%%%%%%%%%
\subsection{Calculus Terminology}
\label{subsec:calculus-terminology-reference}
%%%%%%%%%%%%%%%%%%%%%%%%%%%%%%%%%%%%%%%%%%%%%%%%%%%%%%%%%%%%

\textbf{Derivative.}
The local linear rate of change of a function.

\medskip

\textbf{Differentiable.}
A function for which the derivative exists in the relevant sense.

\medskip

\textbf{Partial derivative.}
A derivative with respect to one variable while holding others fixed.

\medskip

\textbf{Gradient.}
The vector of first partial derivatives of a scalar-valued function.

\medskip

\textbf{Jacobian.}
The matrix of first partial derivatives of a vector-valued function.

\medskip

\textbf{Hessian.}
The matrix of second partial derivatives of a scalar function.

\medskip

\textbf{Antiderivative.}
A function \(F\) satisfying

\[
\boxed{
F'=f.
}
\]

\medskip

\textbf{Integral.}
A mathematical object representing accumulation, area, expectation, mass,
or related quantities depending on context.

\medskip

\textbf{Improper integral.}
An integral involving an infinite interval or unbounded integrand, defined
through limits.

%%%%%%%%%%%%%%%%%%%%%%%%%%%%%%%%%%%%%%%%%%%%%%%%%%%%%%%%%%%%
\subsection{Analysis Terminology}
\label{subsec:analysis-terminology-reference}
%%%%%%%%%%%%%%%%%%%%%%%%%%%%%%%%%%%%%%%%%%%%%%%%%%%%%%%%%%%%

\textbf{Sequence.}
An ordered collection indexed by integers, usually

\[
\boxed{
(a_n).
}
\]

\medskip

\textbf{Series.}
A sum of sequence terms.

\medskip

\textbf{Convergent.}
Approaching a finite limiting object in the relevant topology or norm.

\medskip

\textbf{Divergent.}
Not convergent.

\medskip

\textbf{Cauchy sequence.}
A sequence whose terms become arbitrarily close to one another.

\medskip

\textbf{Complete space.}
A metric space in which every Cauchy sequence converges to a point in the
space.

\medskip

\textbf{Bounded.}
Contained within a finite distance or finite numerical bound appropriate to
the context.

\medskip

\textbf{Supremum.}
The least upper bound.

\medskip

\textbf{Infimum.}
The greatest lower bound.

\medskip

\textbf{Almost everywhere.}
True except on a set of measure zero.

\medskip

\textbf{Measurable.}
Compatible with the relevant sigma-algebra so that measure-theoretic
operations are defined.

%%%%%%%%%%%%%%%%%%%%%%%%%%%%%%%%%%%%%%%%%%%%%%%%%%%%%%%%%%%%
\subsection{Complex Analysis Terminology}
\label{subsec:complex-analysis-terminology}
%%%%%%%%%%%%%%%%%%%%%%%%%%%%%%%%%%%%%%%%%%%%%%%%%%%%%%%%%%%%

\textbf{Holomorphic.}
Complex differentiable on an open set.

\medskip

\textbf{Entire.}
Holomorphic on all of \(\mathbb{C}\).

\medskip

\textbf{Meromorphic.}
Holomorphic except for isolated poles.

\medskip

\textbf{Singularity.}
A point where a complex function fails to be holomorphic.

\medskip

\textbf{Pole.}
An isolated singularity at which the function behaves like a finite negative
power.

\medskip

\textbf{Residue.}
The coefficient of

\[
\boxed{
(z-z_0)^{-1}
}
\]

in the Laurent expansion around \(z_0\).

\medskip

\textbf{Contour.}
An oriented curve in the complex plane used for integration.

%%%%%%%%%%%%%%%%%%%%%%%%%%%%%%%%%%%%%%%%%%%%%%%%%%%%%%%%%%%%
\subsection{Differential Equation Terminology}
\label{subsec:differential-equation-terminology}
%%%%%%%%%%%%%%%%%%%%%%%%%%%%%%%%%%%%%%%%%%%%%%%%%%%%%%%%%%%%

\textbf{Differential equation.}
An equation involving an unknown function and its derivatives.

\medskip

\textbf{Ordinary differential equation.}
A differential equation involving derivatives with respect to one
independent variable.

\medskip

\textbf{Partial differential equation.}
A differential equation involving partial derivatives with respect to
multiple independent variables.

\medskip

\textbf{Order.}
The order of the highest derivative appearing in the equation.

\medskip

\textbf{Linear differential equation.}
An equation linear in the unknown function and its derivatives.

\medskip

\textbf{Initial condition.}
A condition specifying function or derivative values at an initial point.

\medskip

\textbf{Boundary condition.}
A condition imposed on the boundary of a domain.

\medskip

\textbf{Equilibrium.}
A state that remains unchanged under the system dynamics.

\medskip

\textbf{Stable equilibrium.}
An equilibrium for which nearby states remain nearby in the specified
stability sense.

%%%%%%%%%%%%%%%%%%%%%%%%%%%%%%%%%%%%%%%%%%%%%%%%%%%%%%%%%%%%
\subsection{Discrete Mathematics Terminology}
\label{subsec:discrete-mathematics-terminology}
%%%%%%%%%%%%%%%%%%%%%%%%%%%%%%%%%%%%%%%%%%%%%%%%%%%%%%%%%%%%

\textbf{Permutation.}
An ordered arrangement.

\medskip

\textbf{Combination.}
An unordered selection.

\medskip

\textbf{Recurrence relation.}
An equation defining sequence terms in relation to previous terms.

\medskip

\textbf{Generating function.}
A formal or analytic power series encoding a sequence.

\medskip

\textbf{Graph.}
A set of vertices together with edges relating pairs of vertices.

\medskip

\textbf{Vertex.}
A node of a graph.

\medskip

\textbf{Edge.}
A connection between vertices.

\medskip

\textbf{Path.}
A sequence of adjacent vertices or edges.

\medskip

\textbf{Cycle.}
A closed path with appropriate nonrepetition conditions.

\medskip

\textbf{Tree.}
A connected graph with no cycles.

\medskip

\textbf{Matching.}
A set of edges with no shared endpoints.

%%%%%%%%%%%%%%%%%%%%%%%%%%%%%%%%%%%%%%%%%%%%%%%%%%%%%%%%%%%%
\subsection{Probability Terminology}
\label{subsec:probability-terminology-reference}
%%%%%%%%%%%%%%%%%%%%%%%%%%%%%%%%%%%%%%%%%%%%%%%%%%%%%%%%%%%%

\textbf{Experiment.}
A process producing an uncertain outcome.

\medskip

\textbf{Sample space.}
The set of all possible outcomes.

\medskip

\textbf{Event.}
A subset of the sample space.

\medskip

\textbf{Random variable.}
A measurable numerical function on a probability space.

\medskip

\textbf{Probability distribution.}
A mathematical description of how probability is assigned to possible
values.

\medskip

\textbf{Expected value.}
The probability-weighted average or integral

\[
\boxed{
\mathbb{E}[X].
}
\]

\medskip

\textbf{Variance.}
The expected squared deviation from the mean.

\medskip

\textbf{Independent.}
Events \(A,B\) are independent when

\[
\boxed{
\Pr(A\cap B)
=
\Pr(A)\Pr(B).
}
\]

\medskip

\textbf{Conditional probability.}

\[
\boxed{
\Pr(A\mid B)
=
\frac{
\Pr(A\cap B)
}{
\Pr(B)
}
}
\]

when \(\Pr(B)>0\).

%%%%%%%%%%%%%%%%%%%%%%%%%%%%%%%%%%%%%%%%%%%%%%%%%%%%%%%%%%%%
\subsection{Statistics Terminology}
\label{subsec:statistics-terminology-reference}
%%%%%%%%%%%%%%%%%%%%%%%%%%%%%%%%%%%%%%%%%%%%%%%%%%%%%%%%%%%%

\textbf{Population.}
The full collection of units or outcomes under study.

\medskip

\textbf{Sample.}
A subset or collection of observations drawn from a population or data
generating process.

\medskip

\textbf{Parameter.}
A numerical characteristic of a population or statistical model.

\medskip

\textbf{Statistic.}
A numerical function of observed sample data.

\medskip

\textbf{Estimator.}
A rule used to estimate an unknown parameter.

\medskip

\textbf{Estimate.}
The realized numerical value of an estimator.

\medskip

\textbf{Bias.}
The difference between an estimator's expectation and the target parameter.

\[
\boxed{
\operatorname{Bias}(\hat\theta)
=
\mathbb{E}[\hat\theta]-\theta.
}
\]

\medskip

\textbf{Standard error.}
The standard deviation of an estimator or an estimate of that sampling
standard deviation.

\medskip

\textbf{Confidence interval.}
An interval-valued procedure designed to cover a parameter with a specified
long-run frequency under repeated sampling.

\medskip

\textbf{Null hypothesis.}
The hypothesis treated as the reference model in a significance test.

\medskip

\textbf{\(p\)-value.}
A tail probability measuring how incompatible the observed statistic is with
the null model under the specified test.

\medskip

\textbf{Statistical significance.}
A result satisfying a specified inferential threshold, commonly involving
the \(p\)-value and significance level.

%%%%%%%%%%%%%%%%%%%%%%%%%%%%%%%%%%%%%%%%%%%%%%%%%%%%%%%%%%%%
\subsection{Optimization Terminology}
\label{subsec:optimization-terminology-reference}
%%%%%%%%%%%%%%%%%%%%%%%%%%%%%%%%%%%%%%%%%%%%%%%%%%%%%%%%%%%%

\textbf{Objective function.}
The quantity to be minimized or maximized.

\medskip

\textbf{Decision variable.}
A variable whose value is chosen by the optimization process.

\medskip

\textbf{Constraint.}
A condition restricting admissible decisions.

\medskip

\textbf{Feasible point.}
A point satisfying all constraints.

\medskip

\textbf{Feasible region.}
The set of all feasible points.

\medskip

\textbf{Local minimum.}
A point with no smaller objective value in a sufficiently small
neighborhood.

\medskip

\textbf{Global minimum.}
A feasible point whose objective value is no larger than that of any other
feasible point.

\medskip

\textbf{Convex optimization.}
Optimization in which the objective and feasible region possess convex
structure ensuring strong global properties.

\medskip

\textbf{Dual problem.}
A related optimization problem providing bounds and structural information
about the original, or primal, problem.

%%%%%%%%%%%%%%%%%%%%%%%%%%%%%%%%%%%%%%%%%%%%%%%%%%%%%%%%%%%%
\subsection{Numerical Mathematics Terminology}
\label{subsec:numerical-terminology-reference}
%%%%%%%%%%%%%%%%%%%%%%%%%%%%%%%%%%%%%%%%%%%%%%%%%%%%%%%%%%%%

\textbf{Approximation.}
A value or object close to an exact mathematical quantity.

\medskip

\textbf{Absolute error.}

\[
\boxed{
|x-\hat{x}|.
}
\]

\medskip

\textbf{Relative error.}

\[
\boxed{
\frac{
|x-\hat{x}|
}{
|x|
}
}
\]

when \(x\neq0\).

\medskip

\textbf{Truncation error.}
Error introduced by replacing an infinite or exact process with a finite
approximation.

\medskip

\textbf{Roundoff error.}
Error caused by finite-precision arithmetic.

\medskip

\textbf{Conditioning.}
Sensitivity of the exact problem solution to perturbations in the input.

\medskip

\textbf{Stability.}
Sensitivity of a numerical algorithm to perturbations and numerical errors.

\medskip

\textbf{Convergence order.}
The rate at which an approximation approaches the exact solution as a step
size or error parameter tends to zero.

%%%%%%%%%%%%%%%%%%%%%%%%%%%%%%%%%%%%%%%%%%%%%%%%%%%%%%%%%%%%
\subsection{Modeling Terminology}
\label{subsec:modeling-terminology-reference}
%%%%%%%%%%%%%%%%%%%%%%%%%%%%%%%%%%%%%%%%%%%%%%%%%%%%%%%%%%%%

\textbf{Model.}
A mathematical representation of selected features of a real or abstract
system.

\medskip

\textbf{State variable.}
A variable describing the current state of a system.

\medskip

\textbf{Parameter.}
A quantity treated as fixed during a particular model run but potentially
estimated, varied, or calibrated.

\medskip

\textbf{Calibration.}
Adjustment of model parameters to reproduce chosen observations or targets.

\medskip

\textbf{Validation.}
Assessment of how well a model represents the intended real-world behavior
for its purpose.

\medskip

\textbf{Verification.}
Assessment of whether equations or numerical algorithms have been implemented
and solved correctly.

\medskip

\textbf{Sensitivity analysis.}
Study of how model outputs change when inputs or parameters vary.

\medskip

\textbf{Scenario.}
A specified collection of assumptions or parameter conditions under which a
model is evaluated.

\medskip

\textbf{Simulation.}
Numerical or computational generation of system behavior from a mathematical
model.

%%%%%%%%%%%%%%%%%%%%%%%%%%%%%%%%%%%%%%%%%%%%%%%%%%%%%%%%%%%%
\subsection{Common Words With Multiple Mathematical Meanings}
\label{subsec:multiple-meaning-math-terms}
%%%%%%%%%%%%%%%%%%%%%%%%%%%%%%%%%%%%%%%%%%%%%%%%%%%%%%%%%%%%

Several mathematical terms depend strongly on context.

\textbf{Normal} may mean:

\begin{itemize}
    \item perpendicular in geometry;
    \item a normal subgroup in algebra;
    \item a normal matrix in linear algebra;
    \item the normal probability distribution in statistics.
\end{itemize}

\textbf{Order} may mean:

\begin{itemize}
    \item ordering relation;
    \item size of a group;
    \item order of an element;
    \item highest derivative in a differential equation;
    \item asymptotic convergence order;
    \item matrix or tensor order in some contexts.
\end{itemize}

\textbf{Degree} may mean:

\begin{itemize}
    \item polynomial degree;
    \item graph-vertex degree;
    \item field extension degree;
    \item angular degree;
    \item degree of a differential form or map.
\end{itemize}

\textbf{Rank} may mean:

\begin{itemize}
    \item matrix rank;
    \item tensor rank;
    \item rank of an algebraic object;
    \item ordinal or hierarchy position.
\end{itemize}

%%%%%%%%%%%%%%%%%%%%%%%%%%%%%%%%%%%%%%%%%%%%%%%%%%%%%%%%%%%%
\subsection{Words Used in Problem Statements}
\label{subsec:problem-statement-terminology}
%%%%%%%%%%%%%%%%%%%%%%%%%%%%%%%%%%%%%%%%%%%%%%%%%%%%%%%%%%%%

\textbf{Evaluate.}
Find the value of an expression or quantity.

\medskip

\textbf{Simplify.}
Rewrite an expression into an equivalent but typically more compact or useful
form.

\medskip

\textbf{Solve.}
Find all objects satisfying the specified equation or conditions.

\medskip

\textbf{Verify.}
Show that a stated relationship is correct.

\medskip

\textbf{Prove.}
Provide a logically complete mathematical argument.

\medskip

\textbf{Disprove.}
Show that a statement is false, often by counterexample.

\medskip

\textbf{Derive.}
Obtain a formula or result from earlier definitions or principles.

\medskip

\textbf{Approximate.}
Find a value close to the exact value according to a stated method or
accuracy.

\medskip

\textbf{Estimate.}
Use data or mathematical information to infer an unknown quantity.

\medskip

\textbf{Classify.}
Assign objects to categories according to specified mathematical properties.

\medskip

\textbf{Characterize.}
Give necessary and sufficient conditions describing a class of objects.

%%%%%%%%%%%%%%%%%%%%%%%%%%%%%%%%%%%%%%%%%%%%%%%%%%%%%%%%%%%%
\subsection{Common Comparison Terminology}
\label{subsec:comparison-terminology-reference}
%%%%%%%%%%%%%%%%%%%%%%%%%%%%%%%%%%%%%%%%%%%%%%%%%%%%%%%%%%%%

\textbf{Equal.}

\[
\boxed{
a=b.
}
\]

The objects are the same under the relevant notion of equality.

\medskip

\textbf{Equivalent.}
Two objects correspond under a specified equivalence relation but need not be
literally identical.

\medskip

\textbf{Isomorphic.}
Structurally equivalent in an algebraic or categorical sense.

\medskip

\textbf{Homeomorphic.}
Topologically equivalent.

\medskip

\textbf{Congruent.}
Equivalent under a geometric congruence or modular relation, depending on
context.

\medskip

\textbf{Similar.}
Equivalent up to scaling or change of basis, depending on context.

For matrices,

\[
\boxed{
B=P^{-1}AP.
}
\]

%%%%%%%%%%%%%%%%%%%%%%%%%%%%%%%%%%%%%%%%%%%%%%%%%%%%%%%%%%%%
\subsection{Asymptotic Terminology}
\label{subsec:asymptotic-terminology-reference}
%%%%%%%%%%%%%%%%%%%%%%%%%%%%%%%%%%%%%%%%%%%%%%%%%%%%%%%%%%%%

\textbf{Asymptotic.}
Describing behavior in a limiting regime.

\medskip

\textbf{Big-O.}

\[
\boxed{
f(x)=O(g(x))
}
\]

means \(f\) is bounded in magnitude by a constant multiple of \(g\) in the
specified regime.

\medskip

\textbf{Little-o.}

\[
\boxed{
f(x)=o(g(x))
}
\]

means

\[
\boxed{
\frac{f(x)}{g(x)}
\to0.
}
\]

\medskip

\textbf{Asymptotic equivalence.}

\[
\boxed{
f(x)\sim g(x)
}
\]

means

\[
\boxed{
\frac{f(x)}{g(x)}
\to1.
}
\]

%%%%%%%%%%%%%%%%%%%%%%%%%%%%%%%%%%%%%%%%%%%%%%%%%%%%%%%%%%%%
\subsection{Quick Terminology Contrast Table}
\label{subsec:quick-terminology-contrast}
%%%%%%%%%%%%%%%%%%%%%%%%%%%%%%%%%%%%%%%%%%%%%%%%%%%%%%%%%%%%

\begin{center}
\begin{tabular}{c|c}
\textbf{Terms}
&
\textbf{Distinction}
\\ \hline

Domain / Codomain
&
allowed inputs / declared target
\\

Range / Codomain
&
attained outputs / declared target
\\

Injective / Surjective
&
one-to-one / onto
\\

Necessary / Sufficient
&
required / enough
\\

Sequence / Series
&
ordered terms / sum of terms
\\

Parameter / Statistic
&
population-model quantity / sample quantity
\\

Estimate / Estimator
&
realized value / rule
\\

Variance / Standard Deviation
&
squared spread / square-root spread
\\

Standard Deviation / Standard Error
&
data variability / estimator variability
\\

Continuous / Differentiable
&
limit agreement / local linearity
\\

Local / Global Optimum
&
best nearby / best overall
\\

Verification / Validation
&
solve model correctly / represent target appropriately
\end{tabular}
\end{center}

%%%%%%%%%%%%%%%%%%%%%%%%%%%%%%%%%%%%%%%%%%%%%%%%%%%%%%%%%%%%
\subsection{Common Terminology Errors}
\label{subsec:terminology-common-errors}
%%%%%%%%%%%%%%%%%%%%%%%%%%%%%%%%%%%%%%%%%%%%%%%%%%%%%%%%%%%%

\subsubsection{Using ``Equation'' and ``Expression'' Interchangeably}

An expression does not necessarily assert equality.

For example,

\[
x^2+3x+2
\]

is an expression.

\[
x^2+3x+2=0
\]

is an equation.

%%%%%%%%%%%%%%%%%%%%%%%%%%%%%%%%%%%%%%%%%%%%%%%%%%%%%%%%%%%%
\subsubsection{Confusing ``Function'' With ``Formula''}

A function includes its input-output rule together with its relevant domain
and codomain.

Different formulas can describe the same function, and the same formula may
define different functions under different domains.

%%%%%%%%%%%%%%%%%%%%%%%%%%%%%%%%%%%%%%%%%%%%%%%%%%%%%%%%%%%%
\subsubsection{Confusing ``Independent'' With ``Uncorrelated''}

Independence generally implies zero covariance when moments exist.

Zero covariance does not generally imply independence.

%%%%%%%%%%%%%%%%%%%%%%%%%%%%%%%%%%%%%%%%%%%%%%%%%%%%%%%%%%%%
\subsubsection{Confusing ``Random'' With ``Unstructured''}

A random variable or stochastic model may have a precisely specified
probability structure.

``Random'' does not mean mathematically uncontrolled.

%%%%%%%%%%%%%%%%%%%%%%%%%%%%%%%%%%%%%%%%%%%%%%%%%%%%%%%%%%%%
\subsubsection{Confusing ``Proof'' With Numerical Evidence}

Checking many examples can support intuition but does not prove a universal
statement.

%%%%%%%%%%%%%%%%%%%%%%%%%%%%%%%%%%%%%%%%%%%%%%%%%%%%%%%%%%%%
\subsubsection{Using ``Obviously'' Without Justification}

A mathematical step should be justified when its validity is not immediate
from definitions or previously established results.

%%%%%%%%%%%%%%%%%%%%%%%%%%%%%%%%%%%%%%%%%%%%%%%%%%%%%%%%%%%%
\subsection{Exercises}
\label{subsec:mathematical-terminology-exercises}
%%%%%%%%%%%%%%%%%%%%%%%%%%%%%%%%%%%%%%%%%%%%%%%%%%%%%%%%%%%%

\begin{exercise}[1]
Explain the difference between a necessary condition and a sufficient
condition.
\end{exercise}

\begin{exercise}[1]
Explain the difference between a theorem and a conjecture.
\end{exercise}

\begin{exercise}[1]
Explain the difference between a set and an element.
\end{exercise}

\begin{exercise}[1]
Explain the difference between domain, codomain, and range.
\end{exercise}

\begin{exercise}[1]
Explain the difference between injective, surjective, and bijective.
\end{exercise}

\begin{exercise}[1]
Explain the difference between rational and irrational numbers.
\end{exercise}

\begin{exercise}[1]
Explain the difference between an equation and an identity.
\end{exercise}

\begin{exercise}[1]
Explain the difference between linearly independent and spanning.
\end{exercise}

\begin{exercise}[1]
Explain the difference between eigenvalue and eigenvector.
\end{exercise}

\begin{exercise}[1]
Explain the difference between congruent and similar geometric figures.
\end{exercise}

\begin{exercise}[1]
Explain the difference between open and closed sets.
\end{exercise}

\begin{exercise}[1]
Explain the difference between continuity and differentiability.
\end{exercise}

\begin{exercise}[1]
Explain the difference between a sequence and a series.
\end{exercise}

\begin{exercise}[1]
Explain the difference between convergence and boundedness.
\end{exercise}

\begin{exercise}[1]
Explain the difference between an ODE and a PDE.
\end{exercise}

\begin{exercise}[1]
Explain the difference between a permutation and a combination.
\end{exercise}

\begin{exercise}[1]
Explain the difference between probability mass and probability density.
\end{exercise}

\begin{exercise}[1]
Explain the difference between a parameter and a statistic.
\end{exercise}

\begin{exercise}[1]
Explain the difference between an estimator and an estimate.
\end{exercise}

\begin{exercise}[1]
Explain the difference between standard deviation and standard error.
\end{exercise}

\begin{exercise}[2]
Give an example of a function that is injective but not surjective.
\end{exercise}

\begin{exercise}[2]
Give an example of a function that is surjective but not injective.
\end{exercise}

\begin{exercise}[2]
Give an example of a continuous function that is not differentiable at one
point.
\end{exercise}

\begin{exercise}[2]
Give an example of a bounded divergent sequence.
\end{exercise}

\begin{exercise}[2]
Give an example of two uncorrelated random variables that are not independent.
\end{exercise}

\begin{exercise}[2]
Give examples of the word ``normal'' in three different mathematical
contexts.
\end{exercise}

\begin{exercise}[2]
Explain the difference between verification and validation in mathematical
modeling.
\end{exercise}

\begin{exercise}[2]
Explain why statistical significance and practical significance are
different concepts.
\end{exercise}

\begin{exercise}[2]
Explain the difference between local and global minima.
\end{exercise}

\begin{exercise}[2]
Explain the difference between conditioning and numerical stability.
\end{exercise}

%%%%%%%%%%%%%%%%%%%%%%%%%%%%%%%%%%%%%%%%%%%%%%%%%%%%%%%%%%%%
\subsection{Conceptual Summary}
\label{subsec:mathematical-terminology-summary}
%%%%%%%%%%%%%%%%%%%%%%%%%%%%%%%%%%%%%%%%%%%%%%%%%%%%%%%%%%%%

Mathematical language is precise.

Closely related words often carry different technical meanings.

Examples include

\[
\boxed{
\text{necessary}
\neq
\text{sufficient},
}
\]

\[
\boxed{
\text{domain}
\neq
\text{range},
}
\]

\[
\boxed{
\text{sequence}
\neq
\text{series},
}
\]

\[
\boxed{
\text{parameter}
\neq
\text{statistic},
}
\]

and

\[
\boxed{
\text{standard deviation}
\neq
\text{standard error}.
}
\]

A mathematical term should always be interpreted in the context of the
subject and definitions being used.

The central practical principle is

\[
\boxed{
\text{precise terminology supports precise mathematical reasoning}.
}
\]

%%%%%%%%%%%%%%%%%%%%%%%%%%%%%%%%%%%%%%%%%%%%%%%%%%%%%%%%%%%%
\subsection{Section Connections}
\label{subsec:mathematical-terminology-connections}
%%%%%%%%%%%%%%%%%%%%%%%%%%%%%%%%%%%%%%%%%%%%%%%%%%%%%%%%%%%%

\chapterconnection{
Mathematical Terminology consolidates vocabulary used throughout every
chapter of the textbook. The terms are grouped by subject so that readers can
compare related concepts such as necessary versus sufficient, injective
versus surjective, parameter versus statistic, and convergence versus
boundedness. The next section, Glossary, reorganizes the textbook's core
vocabulary into an alphabetical reference, while the final Master Index
provides topic-level navigation across the entire text.
}

%%%%%%%%%%%%%%%%%%%%%%%%%%%%%%%%%%%%%%%%%%%%%%%%%%%%%%%%%%%%
% SECTION SOURCE TRACKING
%%%%%%%%%%%%%%%%%%%%%%%%%%%%%%%%%%%%%%%%%%%%%%%%%%%%%%%%%%%%

%------------------------------------------------------------
% 14.16 Mathematical Terminology
%------------------------------------------------------------
%
% Source basis:
%   - Standard mathematical vocabulary used throughout
%     undergraduate and graduate mathematics.
%
% Used for:
%   - logic and proof terminology;
%   - sets and functions;
%   - number systems;
%   - algebra and linear algebra;
%   - geometry and topology;
%   - calculus and analysis;
%   - complex analysis;
%   - differential equations;
%   - discrete mathematics;
%   - probability and statistics;
%   - optimization;
%   - numerical mathematics;
%   - mathematical modeling;
%   - common problem-statement verbs;
%   - comparison and asymptotic terminology;
%   - common terminology errors.
%
% Notes:
%   - Glossary follows in Section 14.17.
%   - Master Index follows in Section 14.18.
%
%%%%%%%%%%%%%%%%%%%%%%%%%%%%%%%%%%%%%%%%%%%%%%%%%%%%%%%%%%%%